\chapter{2008年陕西卷（文）}

\section{单选题}

\begin{problem}
\(\sin 330^{\circ}\) 等于（\quad）

\choices
    {\(-\frac{\sqrt{3}}{2}\)}
    {\(-\frac{1}{2}\)}
    {\( \frac{1}{2} \)}
    {\(\frac{\sqrt{3}}{2}\)}
\end{problem}

\begin{answer}
B
\end{answer}

\begin{solution}
由三角函数诱导公式，
\(\sin 330^{\circ}=\sin\left(360^{\circ}-30^{\circ}\right)=-\sin30^{\circ}=-\frac{1}{2}\)，故选 B.
\end{solution}

\begin{problem}
已知全集 \(U = \{1, 2, 3, 4, 5\}\)，集合 \(A = \{1, 3\}\)，\(B = \{3, 4, 5\}\)，则集合 \(\complement_U(A \cap B) =\)（\quad）

\choices
    {\(\{3\}\)}
    {\(\{4, 5\}\)}
    {\(\{3, 4, 5\}\)}
    {\(\{1, 2, 4, 5\}\)}
\end{problem}

\begin{answer}
D
\end{answer}

\begin{solution}
先求交集：\(A\cap B=\{3\}\). 因此在全集 \(U\) 中取补集，得到
\(\complement_U(A\cap B)=U\mathbin{\backslash}\{3\}=\{1,2,4,5\}\)，故选 D.
\end{solution}

\begin{problem}
某林场有树苗 \(30000\) 棵，其中松树苗 \(4000\) 棵. 为调查树苗的生长情况，采用分层抽样的方法抽取一个容量为 \(150\) 的样本，则样本中松树苗的数量为（\quad）

\choices
    {\(30\)}
    {\(25\)}
    {\(20\)}
    {\(15\)}
\end{problem}

\begin{answer}
C
\end{answer}

\begin{solution}
分层抽样时，各层被抽取的比例相同. 松树苗在总体中的比例为
\(\frac{4000}{30000}=\frac{2}{15}\)，所以样本中松树苗的数量为
\(150\times\frac{2}{15}=20\)，故选 C.
\end{solution}

\begin{problem}
已知 \(\{a_{n}\}\) 是等差数列，\(a_{1} + a_{2} = 4\)，\(a_{7} + a_{8} = 28\)，则该数列前 \(10\) 项和 \(S_{10}\)（\quad）

\choices
    {\(64\)}
    {\(100\)}
    {\(110\)}
    {\(120\)}
\end{problem}

\begin{answer}
B
\end{answer}

\begin{solution}
设等差数列的公差为 \(d\). 则
\(a_1+a_2=2a_1+d=4\)，且
\(a_7+a_8=(a_1+6d)+(a_1+7d)=2a_1+13d=28\).
两式相减得 \(12d=24\)，所以 \(d=2\)，再由 \(2a_1+d=4\) 得 \(a_1=1\).
于是
\(S_{10}=\frac{10}{2}\left(2a_1+9d\right)=5(2+18)=100\)，故选 B.
\end{solution}

\begin{problem}
直线 \(\sqrt{3} x - y + m = 0\) 与圆 \(x^{2} + y^{2} - 2x - 2 = 0\) 相切，则实数 \(m\) 等于（\quad）

\choices
    {\(-3\sqrt{3}\) 或 \(\sqrt{3}\)}
    {\(-3\sqrt{3}\) 或 \(3\sqrt{3}\)}
    {\(\sqrt{3}\) 或 \(-\sqrt{3}\)}
    {\(-\sqrt{3}\) 或 \(3\sqrt{3}\)}
\end{problem}

\begin{answer}
A
\end{answer}

\begin{solution}
圆的方程可化为
\(\left(x-1\right)^2+y^2=3\)，所以圆心为 \(O(1,0)\)，半径为 \(\sqrt{3}\).
直线与圆相切，当且仅当圆心到直线的距离等于半径，即
\(\frac{|\sqrt{3}+m|}{\sqrt{3+1}}=\sqrt{3}\).
因此 \(|\sqrt{3}+m|=2\sqrt{3}\)，解得
\(m=\sqrt{3}\) 或 \(m=-3\sqrt{3}\)，故选 A.
\end{solution}

\begin{problem}
“\(a=\frac{1}{8}\)”是“对任意的正数 \(x\)，\(2x+\frac{a}{x}\geqslant1\)”的（\quad）

\choices
    {充分不必要条件}
    {必要不充分条件}
    {充要条件}
    {既不充分也不必要条件}
\end{problem}

\begin{answer}
A
\end{answer}

\begin{solution}
当 \(a=\frac{1}{8}\) 且 \(x>0\) 时，由基本不等式，
\[
2x+\frac{a}{x}=2x+\frac{1}{8x}\geqslant2\sqrt{2x\cdot\frac{1}{8x}}=1
\]
所以 \(a=\frac{1}{8}\) 能推出题中的不等式对任意正数 \(x\) 成立.

但反过来不成立. 例如取 \(a=1\)，则
\(2x+\frac{1}{x}\geqslant2\sqrt{2}>1\) 对任意 \(x>0\) 成立，而 \(1\ne\frac{1}{8}\).
因此，\(a=\frac{1}{8}\) 是题中条件的充分不必要条件，故选 A.
\end{solution}

\begin{problem}
已知函数 \( f(x) = 2^{x + 3} \)，\( f^{-1}(x) \) 是 \( f(x) \) 的反函数，若 \(mn = 16\)（\(m\)，\(n \in \R^{+}\)），则 \( f^{-1}(m) + f^{-1}(n) \) 的值为（\quad）

\choices
    {\(10\)}
    {\(4\)}
    {\(1\)}
    {\(-2\)}
\end{problem}

\begin{answer}
D
\end{answer}

\begin{solution}
由 \(y=2^{x+3}\) 得 \(\log_2 y=x+3\)，所以反函数为
\(f^{-1}(x)=\log_2 x-3\)（\(x>0\)）.
因为 \(m\)，\(n\in\R^+\) 且 \(mn=16\)，故
\[
f^{-1}(m)+f^{-1}(n)=\log_2m-3+\log_2n-3
=\log_2(mn)-6=\log_2 16-6=-2
\]
故选 D.
\end{solution}

\begin{problem}
长方体 \(ABCD - A_{1}B_{1}C_{1}D_{1}\) 的各顶点都在半径为 \(1\) 的球面上，其中 \(AB:AD:AA_{1}=2:1:\sqrt{3}\)，则两 \(A\)，\(B\) 点的球面距离为（\quad）

\choices
    {\(\frac{\pi}{4}\)}
    {\(\frac{\pi}{3}\)}
    {\(\frac{\pi}{2}\)}
    {\(\frac{2\pi}{3}\)}
\end{problem}

\begin{answer}
C
\end{answer}

\begin{solution}
设长方体的三条棱长分别为 \(2t\)，\(t\)，\(\sqrt{3}t\). 长方体的体对角线长为
\[
\sqrt{(2t)^2+t^2+(\sqrt{3}t)^2}=2\sqrt{2}t
\]
长方体各顶点都在同一个球面上时，球心是长方体中心，体对角线是球的直径，故
\(2\sqrt{2}t=2\)，从而 \(t=\frac{1}{\sqrt{2}}\)，\(AB=\sqrt{2}\).

设 \(A\)，\(B\) 所对的球心角为 \(\omega\). 球半径为 \(1\)，所以弦长满足
\(AB=2\sin\frac{\omega}{2}\). 因而
\(\sin\frac{\omega}{2}=\frac{\sqrt{2}}{2}\)，且 \(0\leqslant\omega\leqslant\pi\)，得到 \(\omega=\frac{\pi}{2}\).
球面距离等于该弧所对的圆心角，故选 C.
\end{solution}

\begin{problem}
双曲线 \(\frac{x^2}{a^2} - \frac{y^2}{b^2} = 1\)（\(a > 0\)，\(b > 0\)）的左、右焦点分别是 \(F_1\)，\(F_2\)，过 \(F_1\) 作倾斜角为 \(30^\circ\) 的直线交双曲线右支于 \(M\) 点. 若 \(MF_2\) 垂直于 \(x\) 轴，则双曲线的离心率为（\quad）

\choices
    {\(\sqrt{6}\)}
    {\(\sqrt{3}\)}
    {\(\sqrt{2}\)}
    {\(\frac{\sqrt{3}}{3}\)}
\end{problem}

\begin{answer}
B
\end{answer}

\begin{solution}
设双曲线的半焦距为 \(c\)，则焦点为 \(F_1(-c,0)\)，\(F_2(c,0)\)，且
\(c^2=a^2+b^2\). 过 \(F_1\) 的直线倾斜角为 \(30^{\circ}\)，斜率为 \(\frac{1}{\sqrt{3}}\).
由于 \(MF_2\perp x\) 轴，点 \(M\) 的横坐标为 \(c\)，由直线方程得
\(M\left(c,\frac{2c}{\sqrt{3}}\right)\).

将点 \(M\) 代入双曲线方程，并利用 \(b^2=c^2-a^2\)，得
\[
\frac{c^2}{a^2}-\frac{4c^2}{3(c^2-a^2)}=1
\]
令离心率 \(e=\frac{c}{a}>1\)，化为
\(e^2-\frac{4e^2}{3(e^2-1)}=1\)，即
\(3(e^2-1)^2=4e^2\).
令 \(t=e^2\)，则 \(3t^2-10t+3=0\)，解得 \(t=3\) 或 \(t=\frac{1}{3}\). 因为 \(e>1\)，所以 \(e=\sqrt{3}\)，故选 B.
\end{solution}

\begin{problem}
如图，\(\alpha \perp \beta\)，\(\alpha \cap \beta = l\)，\(A \in \alpha\)，\(B \in \beta\)，\(A\)，\(B\) 到 \(l\) 的距离分别是 \(a\) 和 \(b\)，\(AB\) 与 \(\alpha\)，\(\beta\) 所成的角分别是 \(\theta\) 和 \(\varphi\)，\(AB\) 在 \(\alpha\)，\(\beta\) 内的射影分别是 \(m\) 和 \(n\)，若 \(a > b\)，则（\quad）

\FigureLayoutDeclare{img/2008/shaanxi_liberal/q10_fig1.png}{5.3cm}{12bp 7bp 16bp 10bp}
\bitmapfigure[max width=0.34\linewidth]{img/2008/shaanxi_liberal/q10_fig1.png}

\choices
    {\(\theta >\varphi\)，\(m < n\)}
    {\(\theta >\varphi\)，\(m > n\)}
    {\(\theta < \varphi\)，\(m < n\)}
    {\(\theta < \varphi\)，\(m > n\)}
\end{problem}

\begin{answer}
D
\end{answer}

\begin{solution}
建立直角坐标系，使公共直线 \(l\) 为 \(x\) 轴，平面 \(\alpha\) 为 \(xy\) 平面，平面 \(\beta\) 为 \(xz\) 平面. 设 \(A=(x_1,a,0)\)，\(B=(x_2,0,b)\)，并记 \(u=x_2-x_1\)，则
\(AB\) 的长度为 \(L=\sqrt{u^2+a^2+b^2}\).

\(AB\) 与平面 \(\alpha\) 所成角为 \(\theta\)，其垂直于 \(\alpha\) 的分量长度为 \(b\)，所以 \(\sin\theta=\frac{b}{L}\). 同理，\(AB\) 与平面 \(\beta\) 所成角为 \(\varphi\)，故 \(\sin\varphi=\frac{a}{L}\). 由于两角均为锐角且 \(a>b\)，所以 \(\theta<\varphi\).

\(AB\) 在 \(\alpha\)，\(\beta\) 内的射影长度分别为
\(m=\sqrt{u^2+a^2}\)，\(n=\sqrt{u^2+b^2}\). 由 \(a>b\) 得 \(m>n\)，故选 D.
\end{solution}

\begin{problem}
定义在 \(\R\) 上的函数 \(f(x)\) 满足 \(f(x + y) = f(x) + f(y) + 2xy\)（\(x\)，\(y \in \R\)），\(f(1) = 2\)，则 \(f(-2)\) 等于（\quad）

\choices
    {\(2\)}
    {\(3\)}
    {\(6\)}
    {\(9\)}
\end{problem}

\begin{answer}
A
\end{answer}

\begin{solution}
令题中函数关系式中的 \(y=-x\)，则
\(f(0)=f(x)+f(-x)-2x^2\). 再令 \(x=0\)，可得 \(f(0)=0\)，因此
\(f(-x)=2x^2-f(x)\).

由 \(f(1)=2\)，在原函数关系式中取 \(x=y=1\)，得
\(f(2)=f(1)+f(1)+2=6\). 再令 \(x=2\)，\(y=-2\)，得到
\(f(0)=f(2)+f(-2)-8\)，所以
\(f(-2)=2\)，故选 A.
\end{solution}

\begin{problem}
为提高信息在传输中的抗干扰能力，通常在原信息中按一定规则加入相关数据组成传输信息. 设定原信息为 \(a_0 a_1 a_2\)，\(a_i \in \{0, 1\}\)（\(i = 0\)，\(1\)，\(2\)），传输信息为 \(h_0 a_0 a_1 a_2 h_1\)，其中 \(h_0 = a_0 \oplus a_1\)，\(h_1 = h_0 \oplus a_2\)，\(\oplus\) 运算规则为：\(0 \oplus 0 = 0\)，\(0 \oplus 1 = 1\)，\(1 \oplus 0 = 1\)，\(1 \oplus 1 = 0\). 例如原信息为 \(111\)，则传输信息为 \(01111\). 传输信息在传输过程中受到干扰可能导致接收信息出错，则下列接收信息一定有误的是（\quad）

\choices
    {\(11010\)}
    {\(01100\)}
    {\(10111\)}
    {\(00011\)}
\end{problem}

\begin{answer}
C
\end{answer}

\begin{solution}
合法传输信息 \(h_0a_0a_1a_2h_1\) 必须满足
\(h_0=a_0\oplus a_1\) 且 \(h_1=h_0\oplus a_2\). 逐项检查：

接收信息 \(11010\) 中，\(1=1\oplus0\)，且 \(0=1\oplus1\)，合法.
接收信息 \(01100\) 中，\(0=1\oplus1\)，且 \(0=0\oplus0\)，合法.
接收信息 \(10111\) 中，前一校验满足 \(1=0\oplus1\)，但后一校验应为 \(1\oplus1=0\)，而接收到的末位为 \(1\)，故一定有误.
接收信息 \(00011\) 中，\(0=0\oplus0\)，且 \(1=0\oplus1\)，合法.

因此一定有误的是 \(10111\)，故选 C.
\end{solution}

\section{填空题}

\begin{problem}
\(\triangle ABC\) 的内角 \(A\)，\(B\)，\(C\) 的对边分别为 \(a\)，\(b\)，\(c\)，若 \(c = \sqrt{2}\)，\(b = \sqrt{6}\)，\(B = 120^{\circ}\)，则 \(a =\fillinblank{}\).
\end{problem}

\begin{answer}
\(\sqrt{2}\)
\end{answer}

\begin{solution}
由余弦定理，
\[
b^2=a^2+c^2-2ac\cos B
\]
代入 \(b=\sqrt{6}\)，\(c=\sqrt{2}\)，\(B=120^{\circ}\)，得
\(6=a^2+2-2a\sqrt{2}\left(-\frac{1}{2}\right)=a^2+2+\sqrt{2}a\).
所以 \(a^2+\sqrt{2}a-4=0\)，解得
\(a=\sqrt{2}\) 或 \(a=-2\sqrt{2}\). 由于边长 \(a>0\)，故 \(a=\sqrt{2}\).
\end{solution}

\begin{problem}
\(\left(1 - \frac{2}{x}\right)^7\) 的展开式中 \(\frac{1}{x^2}\) 的系数为\fillinblank{}. （用数字作答）
\end{problem}

\begin{answer}
\(84\)
\end{answer}

\begin{solution}
二项式展开的通项为
\(\mathrm C_7^k\left(-\frac{2}{x}\right)^k=\mathrm C_7^k(-2)^k x^{-k}\).
要得到 \(\frac{1}{x^2}=x^{-2}\)，必须取 \(k=2\)，所以所求系数为
\(\mathrm C_7^2(-2)^2=21\times4=84\).
\end{solution}

\begin{problem}
关于平面向量 \(\bs{a}\)，\(\bs{b}\)，\(\bs{c}\)，有下列三个命题：

\begin{circlelist}
\item 若 \(\bs{a} \cdot \bs{b} = \bs{a} \cdot \bs{c}\)，则 \(\bs{b} = \bs{c}\)；
\item 若 \(\bs{a} = (1, k)\)，\(\bs{b} = (-2, 6)\)，\(\bs{a} \parallel \bs{b}\)，则 \(k = -3\)；
\item 非零向量 \(\bs{a}\) 和 \(\bs{b}\) 满足 \( |\bs{a}| = |\bs{b}| = |\bs{a} - \bs{b}| \)，则 \(\bs{a}\) 与 \( \bs{a} + \bs{b} \) 的夹角为 \( 60^{\circ} \).
\end{circlelist}

其中真命题的序号为\fillinblank{}. （写出所有真命题的序号）
\end{problem}

\begin{answer}
\(\circled{2}\)
\end{answer}

\begin{solution}
命题 \(\circled{1}\) 只说明 \(\bs{a}\cdot(\bs{b}-\bs{c})=0\)，不能推出 \(\bs{b}=\bs{c}\). 例如取 \(\bs{a}=(1,0)\)，\(\bs{b}=(0,0)\)，\(\bs{c}=(0,1)\)，则两点积相等而 \(\bs{b}\ne\bs{c}\)，所以命题 \(\circled{1}\) 为假.

命题 \(\circled{2}\) 中，由 \(\bs{a}\parallel\bs{b}\) 得 \(1\cdot6-k\cdot(-2)=0\)，即 \(6+2k=0\)，所以 \(k=-3\)，命题 \(\circled{2}\) 为真.

命题 \(\circled{3}\) 设 \(|\bs{a}|=|\bs{b}|=r>0\). 由 \(|\bs{a}-\bs{b}|=r\)，得
\(r^2=r^2+r^2-2\bs{a}\cdot\bs{b}\)，从而 \(\bs{a}\cdot\bs{b}=\frac{r^2}{2}\). 因此
\[
\cos\angle\left(\bs{a},\bs{a}+\bs{b}\right)
=\frac{\bs{a}\cdot(\bs{a}+\bs{b})}{|\bs{a}|\,|\bs{a}+\bs{b}|}
=\frac{\frac{3r^2}{2}}{r\cdot\sqrt{3}r}=\frac{\sqrt{3}}{2}
\]
所以该夹角为 \(30^{\circ}\)，不是 \(60^{\circ}\)，命题 \(\circled{3}\) 为假. 故真命题的序号为 \(2\).
\end{solution}

\begin{problem}
某地奥运火炬接力传递路线共分 \(6\) 段. 传递活动分别由 \(6\) 名火炬手完成. 如果第一棒火炬手只能从甲，乙，丙三人中产生，最后一棒火炬手只能从甲，乙两人中产生，则不同的传递方案共有\fillinblank{}. （用数字作答）
\end{problem}

\begin{answer}
\(96\)
\end{answer}

\begin{solution}
第一棒从甲、乙、丙中选，最后一棒从甲、乙中选，且六名火炬手分别完成六段，所以第一棒和最后一棒不能由同一人完成. 先安排这两棒：从 \(3\times2=6\) 种搭配中去掉第一棒、最后一棒都是甲或都是乙的两种，得到 \(4\) 种.

确定首尾两人后，剩余四名火炬手安排到中间四段，有 \(4!=24\) 种. 因此不同传递方案共有
\(4\times4!=4\times24=96\) 种.
\end{solution}

\section{解答题}

\begin{problem}
已知函数 \( f(x) = 2\sin \frac{x}{4} \cos \frac{x}{4} + \sqrt{3} \cos \frac{x}{2} \).
\begin{enumerate}
\item 求函数 \( f(x) \) 的最小正周期及最值；
\item 令 \( g(x) = f\left(x + \frac{\pi}{3}\right) \)，判断函数 \( g(x) \) 的奇偶性，并说明理由.
\end{enumerate}
\end{problem}

\begin{solution}
\begin{enumerate}
\item 由正弦、余弦的二倍角公式，得 \(f(x)=\sin\frac{x}{2}+\sqrt{3}\cos\frac{x}{2}=2\sin\left(\frac{x}{2}+\frac{\pi}{3}\right)\). 因为正弦函数的最小正周期为 \(2\pi\)，所以 \(f(x)\) 的最小正周期为 \(4\pi\). 又因为 \(-1\leqslant\sin\left(\frac{x}{2}+\frac{\pi}{3}\right)\leqslant1\)，所以 \(f(x)\) 的最大值为 \(2\)，最小值为 \(-2\).
\item \(g(x)=f\left(x+\frac{\pi}{3}\right)=2\sin\left(\frac{x}{2}+\frac{\pi}{2}\right)=2\cos\frac{x}{2}\). 因为 \(g(-x)=2\cos\left(-\frac{x}{2}\right)=2\cos\frac{x}{2}=g(x)\)，所以 \(g(x)\) 是偶函数.
\end{enumerate}
\end{solution}

\begin{problem}
一个口袋中装有大小相同的 \(2\) 个红球，\(3\) 个黑球和 \(4\) 个白球，从口袋中一次摸出一个球，摸出的球不再放回. 
\begin{enumerate}
\item 连续摸球 \(2\) 次，求第一次摸出黑球，第二次摸出白球的概率；
\item 如果摸出红球，则停止摸球，求摸球次数不超过 \(3\) 次的概率.
\end{enumerate}
\end{problem}

\begin{solution}
\begin{enumerate}
\item 口袋中共有 \(9\) 个球. 第一次摸出黑球的概率为 \(\frac{3}{9}\)，在第一次摸出黑球的条件下，第二次摸出白球的概率为 \(\frac{4}{8}\)，所以所求概率为 \(\frac{3}{9}\cdot\frac{4}{8}=\frac{1}{6}\).
\item 摸球次数不超过 \(3\) 次，等价于前三次摸球中至少摸到一个红球. 也可以按第一次、第二次或第三次首次摸到红球分类：第一次摸到红球的概率为 \(\frac{2}{9}\)；第一次未摸到红球且第二次摸到红球的概率为 \(\frac{7}{9}\cdot\frac{2}{8}\)；前两次均未摸到红球且第三次摸到红球的概率为 \(\frac{7}{9}\cdot\frac{6}{8}\cdot\frac{2}{7}\). 这些事件互不相容，故所求概率为 \(\frac{2}{9}+\frac{7}{9}\cdot\frac{2}{8}+\frac{7}{9}\cdot\frac{6}{8}\cdot\frac{2}{7}=\frac{7}{12}\).
\end{enumerate}
\end{solution}

\begin{problem}
三棱锥被平行于底面 \(ABC\) 的平面所截得的几何体如图所示，截面为 \(A_{1}B_{1}C_{1}\)， \( \angle BAC = 90^{\circ} \)， \( A_{1}A \perp \) 平面 \(ABC\)， \( A_{1}A = \sqrt{3} \)，\(AB = AC = 2A_{1}C_{1} = 2\)，\(D\) 为 \(BC\) 中点.

\begin{enumerate}
\item 证明：平面 \( A_{1}AD \perp \) 平面 \(BCC_{1}B_{1}\)；
\item 求二面角 \(A - CC_{1} - B\) 的大小.
\end{enumerate}

\bitmapfigure[max width=0.34\linewidth]{img/2008/shaanxi_liberal/q19_fig1.png}
\end{problem}

\begin{solution}
\begin{enumerate}
\item 以 \(A\) 为原点，分别以 \(AB\)，\(AC\)，\(AA_{1}\) 的方向为三个坐标轴正方向，建立空间直角坐标系. 由题意，\(A=(0,0,0)\)，\(B=(2,0,0)\)，\(C=(0,2,0)\)，\(A_{1}=(0,0,\sqrt{3})\). 因为截面 \(A_{1}B_{1}C_{1}\) 平行于底面 \(ABC\)，且 \(A_{1}C_{1}=1=\frac{1}{2}AC\)，所以相应线段的长度都缩小为原来的一半，从而 \(B_{1}=(1,0,\sqrt{3})\)，\(C_{1}=(0,1,\sqrt{3})\). 又因为 \(D\) 为 \(BC\) 的中点，所以 \(D=(1,1,0)\).

平面 \(A_{1}AD\) 的方程为 \(x-y=0\)，其一个法向量为 \(\bs{n}_{1}=(1,-1,0)\). 平面 \(BCC_{1}B_{1}\) 的方程为 \(x+y+\frac{z}{\sqrt{3}}=2\)，其一个法向量为 \(\bs{n}_{2}=\left(1,1,\frac{1}{\sqrt{3}}\right)\). 因为 \(\bs{n}_{1}\cdot\bs{n}_{2}=1-1+0=0\)，所以两平面的法向量互相垂直，从而平面 \(A_{1}AD\perp\) 平面 \(BCC_{1}B_{1}\).
\item 设 \(\bs{v}=\overrightarrow{CC_{1}}=(0,-1,\sqrt{3})\). 过 \(C\) 作垂直于 \(CC_{1}\) 的平面，该平面分别与两个面交成二面角的平面角. 取 \(\overrightarrow{CA}=(0,-2,0)\) 和 \(\overrightarrow{CB}=(2,-2,0)\) 在该平面上的投影：
\(\bs{u}=\overrightarrow{CA}-\frac{\overrightarrow{CA}\cdot\bs{v}}{|\bs{v}|^{2}}\bs{v}=\left(0,-\frac{3}{2},-\frac{\sqrt{3}}{2}\right)\)，
\(\bs{w}=\overrightarrow{CB}-\frac{\overrightarrow{CB}\cdot\bs{v}}{|\bs{v}|^{2}}\bs{v}=\left(2,-\frac{3}{2},-\frac{\sqrt{3}}{2}\right)\).
\(\bs{u}\) 和 \(\bs{w}\) 分别沿两个面的平面角方向，因此若该平面角为 \(\theta\)，则
\(\cos\theta=\frac{\bs{u}\cdot\bs{w}}{|\bs{u}||\bs{w}|}=\frac{3}{\sqrt{3}\sqrt{7}}=\sqrt{\frac{3}{7}}\). 因此二面角 \(A-CC_{1}-B\) 的大小为 \(\arccos\sqrt{\frac{3}{7}}\).
\end{enumerate}
\end{solution}

\begin{problem}
已知数列 \(\{a_{n}\}\) 的首项 \(a_1 = \frac{2}{3}\)，\(a_{n + 1} = \frac{2a_n}{a_n + 1}\)，\(n = 1\)，\(2\)，\(\dots\).
\begin{enumerate}
\item 证明：数列 \( \left\{\frac{1}{a_{n}}-1\right\} \) 是等比数列；
\item 数列 \( \left\{\frac{n}{a_{n}}\right\} \) 的前 \(n\) 项和 \( S_{n} \).
\end{enumerate}
\end{problem}

\begin{solution}
\begin{enumerate}
\item 由递推关系式，得 \(\frac{1}{a_{n+1}}=\frac{a_n+1}{2a_n}=\frac{1}{2}\left(\frac{1}{a_n}+1\right)\)，所以 \(\frac{1}{a_{n+1}}-1=\frac{1}{2}\left(\frac{1}{a_n}-1\right)\). 又因为 \(\frac{1}{a_1}-1=\frac{1}{2}\)，所以数列 \(\left\{\frac{1}{a_n}-1\right\}\) 是首项为 \(\frac{1}{2}\)、公比为 \(\frac{1}{2}\) 的等比数列.
\item 由上问，\(\frac{1}{a_n}-1=\frac{1}{2^n}\)，从而 \(\frac{1}{a_n}=1+\frac{1}{2^n}\)，即 \(a_n=\frac{2^n}{2^n+1}\). 因此 \(\frac{n}{a_n}=n+\frac{n}{2^n}\)，于是
\[
S_n=\sum_{k=1}^{n}\frac{k}{a_k}=\frac{n(n+1)}{2}+\sum_{k=1}^{n}\frac{k}{2^k}
\]
令 \(T_n=\sum_{k=1}^{n}\frac{k}{2^k}\)，注意到 \(\frac{k}{2^k}=\frac{k+1}{2^{k-1}}-\frac{k+2}{2^k}\)，所以
\[
T_n=2-\frac{n+2}{2^n}
\]
故 \(S_n=\frac{n(n+1)}{2}+2-\frac{n+2}{2^n}\).
\end{enumerate}
\end{solution}

\begin{problem}
已知抛物线 \(C: y = 2x^{2}\)，直线 \(y = kx + 2\) 交 \(C\) 于 \(A\)，\(B\) 两点. \(M\) 是线段 \(AB\) 的中点，过 \(M\) 作 \(x\) 轴的垂线交 \(C\) 于点 \(N\). 
\begin{enumerate}
\item 证明：抛物线 \(C\) 在点 \(N\) 处的切线与 \(AB\) 平行；
\item 是否存在实数 \(k\) 使 \(\overrightarrow{NA} \cdot \overrightarrow{NB} = 0\). 若存在，求 \(k\) 的值；若不存在，说明理由.
\end{enumerate}
\end{problem}

\begin{solution}
\begin{enumerate}
\item 设 \(A\)，\(B\) 的横坐标分别为 \(x_A\)，\(x_B\). 由交点横坐标满足的方程 \(2x^2-kx-2=0\)，得 \(x_A+x_B=\frac{k}{2}\). 因为 \(M\) 是 \(AB\) 的中点，所以 \(M\) 的横坐标为 \(\frac{x_A+x_B}{2}=\frac{k}{4}\)，过 \(M\) 作 \(x\) 轴的垂线与抛物线交于 \(N\)，故 \(N\) 的横坐标为 \(\frac{k}{4}\). 抛物线 \(y=2x^2\) 在横坐标为 \(x\) 的点处的切线斜率为 \(4x\)，所以点 \(N\) 处切线的斜率为 \(4\cdot\frac{k}{4}=k\). 直线 \(AB\) 的斜率也是 \(k\)，因此抛物线在点 \(N\) 处的切线与 \(AB\) 平行.
\item 令 \(p=\frac{x_A+x_B}{2}=\frac{k}{4}\)，\(d=\frac{x_A-x_B}{2}\)，则 \(x_A=p+d\)，\(x_B=p-d\). 由韦达定理，\(x_Ax_B=-1\)，所以 \(p^2-d^2=-1\)，即 \(d^2=p^2+1\). 又 \(N=(p,2p^2)\)，从而

\[
\overrightarrow{NA}=\left(d,4pd+2d^2\right)
\]

\[
\overrightarrow{NB}=\left(-d,-4pd+2d^2\right)
\]
因此
\[
\overrightarrow{NA}\cdot\overrightarrow{NB}=-d^2+\left(4pd+2d^2\right)\left(-4pd+2d^2\right)=d^2\left(4d^2-16p^2-1\right)=3d^2\left(1-4p^2\right)
\]
因为 \(A\)，\(B\) 为两个不同的交点，所以 \(d\neq0\). 当且仅当 \(p^2=\frac{1}{4}\) 时，上式等于 \(0\)，此时 \(p=\pm\frac{1}{2}\)，所以 \(k=4p=\pm2\). 对这两个值，交点方程的判别式均为正，确有两个交点，故存在满足条件的实数 \(k\)，其值为 \(2\) 或 \(-2\).
\end{enumerate}
\end{solution}

\begin{problem}
设函数 \( f(x) = x^{3} + ax^{2} - a^{2}x + 1 \)，\( g(x) = ax^{2} - 2x + 1 \)，其中实数 \( a \neq 0 \). 
\begin{enumerate}
\item 若 \(a > 0\)，求函数 \(f(x)\) 的单调区间；
\item 当函数 \( y = f(x) \) 与 \( y = g(x) \) 的图象只有一个公共点且 \( g(x) \) 存在最小值时，记 \( g(x) \) 的最小值为 \( h(a) \)，求 \( h(a) \) 的值域；
\item 若 \( f(x) \) 与 \( g(x) \) 在区间 \((a, a + 2)\) 内均为增函数，求 \( a \) 的取值范围.
\end{enumerate}
\end{problem}

\begin{solution}
\begin{enumerate}
\item 当 \(a>0\) 时，\(f'(x)=3x^2+2ax-a^2=(3x-a)(x+a)\). 导数的两个零点为 \(-a\) 和 \(\frac{a}{3}\)，且 \(-a<\frac{a}{3}\). 因此当 \(x<-a\) 或 \(x>\frac{a}{3}\) 时，\(f'(x)>0\)；当 \(-a<x<\frac{a}{3}\) 时，\(f'(x)<0\). 所以 \(f(x)\) 的增区间为 \((-\infty,-a)\) 和 \(\left(\frac{a}{3},+\infty\right)\)，减区间为 \(\left(-a,\frac{a}{3}\right)\).
\item 两函数图象的公共点横坐标满足 \(f(x)=g(x)\)，即 \(x^3+ax^2-a^2x+1=ax^2-2x+1\)，整理得 \(x\left(x^2+2-a^2\right)=0\). 因而 \(x=0\) 总是公共点横坐标；当 \(a^2>2\) 时，方程 \(x^2=a^2-2\) 还给出两个另外的公共点横坐标.

因为 \(g(x)=ax^2-2x+1\) 存在最小值，所以必须有 \(a>0\). 要使公共点只有一个，必须且只需 \(a^2\leqslant2\)，故 \(0<a\leqslant\sqrt{2}\). 配方得 \(g(x)=a\left(x-\frac{1}{a}\right)^2+1-\frac{1}{a}\)，所以 \(h(a)=1-\frac{1}{a}\). 当 \(0<a\leqslant\sqrt{2}\) 时，\(h(a)\) 在该区间上递增，且当 \(a\to0^+\) 时趋于 \(-\infty\)，在 \(a=\sqrt{2}\) 时取最大值 \(1-\frac{\sqrt{2}}{2}\). 因此 \(h(a)\) 的值域为 \(\left(-\infty,1-\frac{\sqrt{2}}{2}\right]\).
\item \(f'(x)=(3x-a)(x+a)\)，\(g'(x)=2ax-2=2(ax-1)\). 先考虑 \(a>0\). 由于 \(x\in(a,a+2)\) 时有 \(x>a>\frac{a}{3}\)，所以 \(f'(x)>0\). 要使 \(g'(x)>0\) 对区间内一切 \(x\) 成立，需要 \(a\geqslant1\)；当 \(a=1\) 时，区间为开区间 \((1,3)\)，仍有 \(g'(x)>0\). 故正数情形给出 \(a\geqslant1\).

再考虑 \(a<0\). 此时 \(f'(x)>0\) 的区域为 \(x<\frac{a}{3}\) 或 \(x>-a\). 由于 \(a<0\) 时不可能有整个区间 \((a,a+2)\) 位于 \(x>-a\) 的区域内，要使其位于左侧增区间内，必须有 \(a+2\leqslant\frac{a}{3}\)，即 \(a\leqslant-3\). 对于 \(a\leqslant-3\)，又有 \(a(a+2)\geqslant1\)，从而对 \(x\in(a,a+2)\)，有 \(ax>a(a+2)\geqslant1\)，所以 \(g'(x)>0\). 反之，若 \(-3<a<0\)，则区间内存在 \(x\) 使 \(\frac{a}{3}<x<-a\)，此时 \(f'(x)<0\)，不满足题意. 因此负数情形给出 \(a\leqslant-3\).

综上，\(a\) 的取值范围为 \(\left(-\infty,-3\right]\cup[1,+\infty)\).
\end{enumerate}
\end{solution}
